\documentclass[11pt]{article}
\usepackage[margin=0.9in]{geometry}
\usepackage{amsmath,amssymb,amsthm,booktabs,mathrsfs}
\usepackage[T1]{fontenc}
\usepackage{lmodern,microtype}
\usepackage[colorlinks=true,linkcolor=blue,urlcolor=blue]{hyperref}
\newcommand{\norm}[1]{\left\lVert#1\right\rVert}
\newtheorem{theorem}{Theorem}[section]
\newtheorem{lemma}[theorem]{Lemma}
\newtheorem{proposition}[theorem]{Proposition}
\newtheorem{corollary}[theorem]{Corollary}
\theoremstyle{remark}
\newtheorem{remark}[theorem]{Remark}
\title{Toward a four-dimensional Yang--Mills construction\\
\large Notebook 35: infrared cancellation and physical heat bounds}
\author{Research working document}
\date{17 October 2026}
\begin{document}
\maketitle
\begin{abstract}
We derive an exact local, gauge-covariant representation
of the actual leading determinant using zero covariant
block means in the scalar gauge. Both its scalar and
vector factors have flat spectral gaps of order
\(K^{-2}\), with uniform heat and inverse-power bounds.
The finite amplitude identity continues jointly through
the flat stabilizers and passes to local infinite
systems at each fixed \(K\). Its invariant diagonal
densities give the actual Maxwell moment directly.
This closes the specific infrared passage left open
by the earlier massless physical quotient.

The exact finite physical quotient of Notebook 34 requires
joint control of its scalar and physical determinants.
We give an explicit reason: a globally based massless
scalar determinant on a periodic constant Cartan connection
has a nonzero quadratic response per volume, whereas the
ordered thermodynamic and zero-mass limit gives zero.
The fine/coarse scalar ratio has an extra term proportional
to \(K^2-1\). The actual local gauge Ward identity and
the finite quotient formula prove that the physical
determinant cancels this term exactly in the thermodynamic
limit. No flux seam or nonperiodic magnetic family is used.

We also extend the flat physical spectral estimates to
uniform heat-trace and inverse-power bounds. The
rescaled inverse square has at most logarithmic trace
growth; powers above two have uniformly bounded trace
per final block. A local tent lift proves full
Laplacian domination for the new vector gauge, so
both new factors have the corresponding estimates.
Their endpoint-uniform curvature comparison remains
the next step toward a positive actual drift.
\end{abstract}

\section*{Scope and conventions}

All actual operators are those of the homogeneous local
construction in Notebooks 32--34. The primitive blocking
ratio is fixed, \(K=\ell^k\), and \(M\) denotes the
preceding homogeneous depth. Traces count real adjoint
coordinates. For a normalized Cartan element \(H\),
\(\tau(H^2)=1\), set
\[
 C_H=\sum_{q\in\operatorname{spec}(\operatorname{ad}_H)}q^2,
 \qquad \alpha_g=\frac{C_H}{48\pi^2}.
\]
The complexified charge decomposition counts each real
adjoint degree of freedom once in this trace convention.
The notation \([t^2]\) denotes a Taylor coefficient,
not a second derivative.

The constant flat connection used below tests the
zero-momentum Ward identity. It is different from the
constant-curvature probe that extracts the Maxwell
coefficient. The infrared order of limits is stated
explicitly for each result.

